\documentclass[11pt,a4paper]{article}

\usepackage[utf8]{inputenc}
\usepackage{amsmath,amssymb}
\usepackage{booktabs,tabularx}
\usepackage{enumitem}
\usepackage{geometry}
\usepackage{hyperref}
\usepackage{xcolor}
\usepackage{listings}
\usepackage{fancyhdr}
\usepackage{titlesec}

\geometry{margin=2.3cm,top=2.5cm,bottom=2.5cm}

\hypersetup{
  colorlinks=true,
  linkcolor=blue!60!black,
  urlcolor=blue!60!black,
  pdftitle={ISLS v3.1 --- Template-Free Generation},
  pdfauthor={Sebastian Klemm},
}

\lstset{
  basicstyle=\ttfamily\scriptsize,
  breaklines=true,
  frame=single,
  backgroundcolor=\color{gray!8},
  keywordstyle=\color{blue!70!black}\bfseries,
  commentstyle=\color{green!50!black}\itshape,
  numbers=left,
  numberstyle=\tiny\color{gray},
  numbersep=5pt,
  tabsize=4,
  showstringspaces=false,
}

\pagestyle{fancy}
\fancyhf{}
\fancyhead[L]{\small ISLS v3.1 --- Template-Free Generation}
\fancyhead[R]{\small Klemm, 2026}
\fancyfoot[C]{\thepage}

\begin{document}

\begin{titlepage}
\begin{center}
\vspace*{1.5cm}
{\Huge\bfseries ISLS v3.1}\\[0.4cm]
{\LARGE\bfseries Template-Free Generation}\\[0.3cm]
{\LARGE\bfseries Norms + LLM + Compile Verification}\\[0.3cm]
{\LARGE\bfseries One Domain Perfect End-to-End}\\[1.2cm]
{\Large Specification v1.0}\\[1.5cm]
{\large Sebastian Klemm}\\[0.3cm]
{\small March 2026}\\[2cm]

\begin{tabular}{ll}
\textbf{Repository:} & \texttt{graphity/isls} \\
\textbf{Replaces:} & Tera templates in isls-orchestrator \\
\textbf{Core change:} & Every file is LLM-generated, \\
 & norm-constrained, compile-verified. \\
\textbf{Goal:} & \texttt{docker-compose up} $\to$ working app in browser \\
\textbf{Requires:} & OpenAI API key \\
\end{tabular}

\vfill
{\small
Templates are dead. They break on every domain.\\
Norms define WHAT. The LLM generates HOW. The compiler verifies.\\
One domain. Warehouse. Perfect. End-to-end.\\
Login $\to$ create product $\to$ adjust stock $\to$ create order $\to$
fulfill $\to$ report. In the browser.
}
\end{center}
\end{titlepage}

\tableofcontents
\newpage

% ============================================================
\section{Why Templates Must Die}

Templates are rigid. Every domain breaks them differently:

\begin{center}
\small
\begin{tabular}{llp{5.5cm}}
\toprule
\textbf{Domain} & \textbf{Template Bug} & \textbf{Root Cause} \\
\midrule
Warehouse & \texttt{CreateUser} has no \texttt{last\_login\_at} & Quality hardening removed field, template still references it \\
E-Commerce & No \texttt{User} entity at all & Domain has \texttt{Customer}, auth template hardcodes \texttt{User} \\
Project & \texttt{User} has \texttt{display\_name} not \texttt{name} & Template hardcodes \texttt{user.name} \\
All & \texttt{as\_deref()} on \texttt{Option<i32>} & Template doesn't know which fields are \texttt{String} vs \texttt{i32} \\
All & Service bodies are CRUD boilerplate & Templates can't express domain logic \\
All & Frontend shows empty tables & Templates have no data-aware rendering \\
\bottomrule
\end{tabular}
\end{center}

Every fix creates a new break. The template system cannot adapt to
domain variation because it is fundamentally static.

\subsection{The New Paradigm}

\begin{lstlisting}[title=Old: Template-Driven]
Norm says: "Generate a CRUD service for Product"
Template says: "Here's a fixed Rust file with {entity} placeholders"
Result: Code that compiles if the entity matches the template's assumptions.
        Breaks if it doesn't.
\end{lstlisting}

\begin{lstlisting}[title=New: LLM-Driven with Norm Constraints]
Norm says: "Generate a CRUD service for Product with these 5 functions,
           these exact types, these error handling patterns"
LLM says:  "Here's the code, using the exact types you specified"
Compiler says: "Compiles." or "Error on line 47: no field 'price'"
If error: LLM gets the error, fixes it, resubmits.
After 3 attempts: compiled, verified, accepted.
\end{lstlisting}

The LLM replaces the template. The norm replaces the template's
structural assumptions. The compiler replaces hope.

% ============================================================
\section{Architecture}

\subsection{Generation Pipeline (Revised)}

\begin{lstlisting}[title=v3.1 Pipeline]
TOML Input
  | isls-hypercube: parse, domain detection, entity inference
  v
HyperCube (dimensions, couplings)
  | isls-norms: match norms, compose, wire dependencies
  v
ComposedPlan (every artifact defined: models, queries, services, api, frontend, tests)
  | isls-decomposer: Fiedler collapse order
  v
Ordered Generation Steps
  | For each step:
  |   1. Build prompt from norm artifact spec + type context of already-generated files
  |   2. Call LLM (OpenAI API)
  |   3. Strip markdown fences
  |   4. Write to file
  |   5. Run cargo check (backend) or syntax check (frontend)
  |   6. If error: feed error back to LLM, retry (max 3)
  |   7. If success: update type context (new types now available for next file)
  |   8. Crystallize pattern
  v
Complete Application (backend + frontend + docker + migrations)
  | Final: cargo build && docker-compose config --quiet
  v
Verified Output
\end{lstlisting}

\subsection{Key Insight: Sequential Generation with Growing Context}

Files are generated in dependency order (determined by Fiedler
decomposition). Each generated file becomes part of the type
context for the next file:

\begin{lstlisting}
Step 1: Generate src/errors.rs
        Context: norm spec only
        -> errors.rs exists, AppError is now known

Step 2: Generate src/pagination.rs
        Context: norm spec + errors.rs types
        -> PaginationParams, PaginatedResponse now known

Step 3: Generate src/auth.rs
        Context: norm spec + errors.rs + pagination.rs
        -> AuthUser, Claims, require_role now known

Step 4: Generate src/models/product.rs
        Context: all above + Product entity spec from norm
        -> Product, CreateProduct, UpdateProduct now known

Step 5: Generate src/database/product_queries.rs
        Context: all above (knows exact Product fields)
        -> list, get_by_id, insert, update, delete now known

Step 6: Generate src/services/product.rs
        Context: all above (knows queries + models + errors + auth)
        -> Service functions with REAL business logic

Step 7: Generate src/api/product.rs
        Context: all above (knows services + auth + pagination)
        -> Endpoints that call real service functions
\end{lstlisting}

No file is generated in isolation. Every file sees what came before.
The LLM cannot hallucinate types because the prompt contains the
exact definitions from previously generated files.

% ============================================================
\section{The Generation Engine}

\subsection{New Module: isls-forge-llm}

Replace the Tera template emission in isls-orchestrator with an
LLM-driven generation engine. This is a new module inside
\texttt{isls-renderloop} or a new crate \texttt{isls-forge-llm}.

\begin{lstlisting}[language=C]
pub struct LlmForge {
    oracle: Box<dyn Oracle>,
    type_context: TypeContext,
    norm_plan: ComposedPlan,
    output_dir: PathBuf,
    generated_files: Vec<GeneratedFile>,
    stats: ForgeStats,
}

pub struct GeneratedFile {
    pub path: String,
    pub content: String,
    pub generation_method: String, // "llm" or "static" (for Cargo.toml, etc.)
    pub attempts: u32,
    pub tokens_used: u64,
}

pub struct ForgeStats {
    pub files_generated: usize,
    pub total_tokens: u64,
    pub compile_checks: usize,
    pub compile_failures: usize,
    pub retries: usize,
    pub total_time_secs: f64,
}

impl LlmForge {
    /// Generate the entire application, file by file, in dependency order.
    pub fn generate(&mut self, spec: &AppSpec) -> Result<Vec<GeneratedFile>> {
        // 1. Generate static files (no LLM needed)
        self.generate_static_files(spec)?;

        // 2. Determine file generation order from Fiedler decomposition
        let order = self.generation_order();

        // 3. Generate each file
        for file_spec in &order {
            let result = self.generate_file(file_spec)?;
            // After each file: update type context
            self.type_context.add_file(&result.path, &result.content);
            self.generated_files.push(result);
        }

        // 4. Final compilation check
        self.final_compile_check()?;

        Ok(self.generated_files.clone())
    }
}
\end{lstlisting}

\subsection{File Generation Order}

Based on dependency topology (Fiedler decomposition collapses
most-depended-on first):

\begin{lstlisting}[title=Generation Order for a Typical App]
LAYER 0 — Infrastructure (static, no LLM):
  Cargo.toml, .env.example, docker-compose.yml, Dockerfile, README.md

LAYER 1 — Foundation (LLM, no dependencies):
  src/errors.rs        (AppError enum, ResponseError impl)
  src/config.rs        (AppConfig, env loading)
  src/pagination.rs    (PaginationParams, PaginatedResponse)

LAYER 2 — Auth (LLM, depends on Layer 1):
  src/models/user.rs   (User struct — ADAPTS to domain's user entity)
  src/auth.rs          (JWT, Claims, AuthUser, middleware, role guard)

LAYER 3 — Models (LLM, depends on Layer 1):
  src/models/{entity}.rs  for each entity (struct, Create, Update, Validate)
  src/models/mod.rs       (module declarations)

LAYER 4 — Database (LLM, depends on Layer 3):
  database/migrations/001_initial.sql  (all tables, FK, indices)
  src/database/{entity}_queries.rs     (CRUD queries using exact model types)
  src/database/pool.rs                 (connection pool)
  src/database/mod.rs

LAYER 5 — Services (LLM, depends on Layer 3+4):
  src/services/{entity}.rs  (business logic, uses exact query fns + model types)
  src/services/mod.rs

LAYER 6 — API (LLM, depends on Layer 2+5):
  src/api/{entity}.rs       (endpoints, uses exact service fns + auth)
  src/api/auth_routes.rs    (login, register — uses exact User type)
  src/api/mod.rs

LAYER 7 — Main (LLM, depends on all):
  src/main.rs               (wires everything together)

LAYER 8 — Frontend (LLM, depends on API spec):
  frontend/index.html
  frontend/style.css
  frontend/src/api/client.js
  frontend/src/pages/{module}.js
  frontend/src/components/*.js

LAYER 9 — Tests (LLM, depends on all):
  tests/api_tests.rs
\end{lstlisting}

\subsection{Static Files (No LLM)}

Some files are always the same structure. Generate them directly:

\begin{lstlisting}[language=C]
fn generate_static_files(&mut self, spec: &AppSpec) -> Result<()> {
    // Cargo.toml — always same structure, deps from norm
    self.write_file("backend/Cargo.toml", &generate_cargo_toml(spec));
    // docker-compose.yml
    self.write_file("docker-compose.yml", &generate_docker_compose(spec));
    // Dockerfile
    self.write_file("backend/Dockerfile", &DOCKERFILE_TEMPLATE);
    // .env.example
    self.write_file(".env.example", &generate_env_example(spec));
    // .gitignore
    self.write_file(".gitignore", &GITIGNORE_TEMPLATE);
    Ok(())
}
\end{lstlisting}

These are still ``templates'' but trivial ones (10-30 lines each)
that never break because they have no entity-specific content.

\subsection{LLM File Generation}

\begin{lstlisting}[language=C]
fn generate_file(&mut self, file_spec: &FileSpec) -> Result<GeneratedFile> {
    let prompt = self.build_prompt(file_spec);

    for attempt in 1..=3 {
        let response = self.oracle.call(&prompt, 4096)?;
        let code = strip_markdown_fences(&response);

        // Write to disk
        let path = self.output_dir.join(&file_spec.path);
        std::fs::create_dir_all(path.parent().unwrap())?;
        std::fs::write(&path, &code)?;

        // Compile check (for Rust files)
        if file_spec.path.ends_with(".rs") {
            match self.cargo_check() {
                Ok(()) => {
                    return Ok(GeneratedFile {
                        path: file_spec.path.clone(),
                        content: code,
                        generation_method: "llm".into(),
                        attempts: attempt,
                        tokens_used: estimate_tokens(&prompt, &response),
                    });
                }
                Err(compile_error) => {
                    if attempt < 3 {
                        // Feed error back to LLM
                        let fix_prompt = format!(
                            "{}\n\n\
                             ## COMPILATION ERROR\n\
                             The code above produced this error:\n```\n{}\n```\n\n\
                             Fix the error. Return the COMPLETE corrected file.\n\
                             Do not change function signatures or types from the context above.\n\
                             Do not add new dependencies.",
                            prompt, compile_error
                        );
                        // Retry with error context
                        let retry = self.oracle.call(&fix_prompt, 4096)?;
                        let fixed = strip_markdown_fences(&retry);
                        std::fs::write(&path, &fixed)?;
                        // Will be checked on next loop iteration
                    }
                }
            }
        } else {
            // Non-Rust files: accept without compile check
            return Ok(GeneratedFile { ... });
        }
    }

    // All 3 attempts failed: return last attempt anyway (best effort)
    Err(anyhow!("Failed to generate {} after 3 attempts", file_spec.path))
}
\end{lstlisting}

\subsection{cargo check Integration}

\begin{lstlisting}[language=C]
/// Run cargo check on the generated backend.
/// Returns Ok(()) if it compiles, Err(error_message) if not.
fn cargo_check(&self) -> Result<(), String> {
    let output = std::process::Command::new("cargo")
        .arg("check")
        .arg("--message-format=short")
        .current_dir(self.output_dir.join("backend"))
        .output()
        .map_err(|e| format!("Failed to run cargo: {}", e))?;

    if output.status.success() {
        Ok(())
    } else {
        let stderr = String::from_utf8_lossy(&output.stderr);
        // Extract just the error lines, not warnings
        let errors: String = stderr.lines()
            .filter(|l| l.contains("error[E") || l.contains("error:"))
            .take(10)  // max 10 error lines to keep prompt short
            .collect::<Vec<_>>()
            .join("\n");
        Err(errors)
    }
}
\end{lstlisting}

% ============================================================
\section{Prompt Construction}

This is the critical part. The prompt must give the LLM everything
it needs to produce correct, compilable code on the first attempt.

\subsection{Prompt Structure}

Every prompt follows this structure:

\begin{lstlisting}[title=Universal Prompt Template]
You are generating one file for a {app_name} application in Rust.

## APPLICATION CONTEXT
{app_description}

## ALREADY GENERATED FILES — USE THESE EXACT TYPES
{type_context}

## NORM REQUIREMENTS FOR THIS FILE
{norm_artifacts_for_this_file}

## FILE TO GENERATE
Path: {file_path}
Purpose: {file_purpose}

## RULES
- Use ONLY types shown in "ALREADY GENERATED FILES"
- Use `tracing::{info, warn, error}`, never `log::{}`
- Use `sqlx::query_as::<_, Type>()` (runtime), never `sqlx::query_as!()`
- All errors via AppError variants shown above
- All async functions return Result<T, AppError>
- No unwrap() — use ? operator
- Do not add use statements for external crates not in Cargo.toml
- Return the COMPLETE file content. No markdown fences.
  No explanation. Just Rust code.

## DOMAIN-SPECIFIC REQUIREMENTS
{business_rules_for_this_entity}

Generate the complete file now.
\end{lstlisting}

\subsection{Type Context: Growing with Each File}

\begin{lstlisting}[language=C]
impl TypeContext {
    /// After a file is generated and compiled, add its types to the context.
    pub fn add_file(&mut self, path: &str, content: &str) {
        // Extract public structs, enums, functions from content
        // Add them to the context so the NEXT file can reference them
        if path.contains("errors") {
            self.error_types = extract_public_items(content);
        } else if path.contains("pagination") {
            self.pagination_types = extract_public_items(content);
        } else if path.contains("auth") {
            self.auth_types = extract_public_items(content);
        } else if path.contains("models/") {
            let entity = entity_from_path(path);
            self.models.insert(entity, extract_public_items(content));
        } else if path.contains("database/") && path.contains("queries") {
            let entity = entity_from_path(path);
            self.queries.insert(entity, extract_fn_signatures(content));
        } else if path.contains("services/") {
            let entity = entity_from_path(path);
            self.services.insert(entity, extract_fn_signatures(content));
        }
    }

    /// Build the type context string for a prompt.
    /// Includes ALL types generated so far.
    pub fn to_prompt_string(&self) -> String {
        let mut s = String::new();
        if !self.error_types.is_empty() {
            s.push_str("### Error Types (from src/errors.rs)\n```rust\n");
            s.push_str(&self.error_types);
            s.push_str("\n```\n\n");
        }
        if !self.pagination_types.is_empty() {
            s.push_str("### Pagination (from src/pagination.rs)\n```rust\n");
            s.push_str(&self.pagination_types);
            s.push_str("\n```\n\n");
        }
        // ... auth, models, queries, services
        for (entity, types) in &self.models {
            s.push_str(&format!("### Model: {} (from src/models/{}.rs)\n```rust\n",
                entity, entity));
            s.push_str(types);
            s.push_str("\n```\n\n");
        }
        for (entity, sigs) in &self.queries {
            s.push_str(&format!("### Queries: {} (from src/database/{}_queries.rs)\n```rust\n",
                entity, entity));
            s.push_str(&sigs.join("\n"));
            s.push_str("\n```\n\n");
        }
        for (entity, sigs) in &self.services {
            s.push_str(&format!("### Services: {} (from src/services/{}.rs)\n```rust\n",
                entity, entity));
            s.push_str(&sigs.join("\n"));
            s.push_str("\n```\n\n");
        }
        s
    }
}
\end{lstlisting}

\subsection{Norm Artifacts in Prompt}

The norm tells the LLM what the file must contain:

\begin{lstlisting}[language=C]
fn norm_context_for_file(&self, file_spec: &FileSpec) -> String {
    match file_spec.file_type {
        FileType::Model => format!(
            "This file must define:\n\
             - pub struct {} with these fields: {}\n\
             - pub struct Create{} (no id, no timestamps)\n\
             - pub struct Update{} (all fields Option<T>)\n\
             - impl {} {{ pub fn validate(&self) -> Result<(), Vec<String>> }}\n\
             - All structs derive: Debug, Clone, Serialize, Deserialize, sqlx::FromRow\n",
            entity.name, fields_description, entity.name, entity.name, entity.name
        ),
        FileType::Queries => format!(
            "This file must define these async functions:\n\
             - pub async fn list(pool, params, filters...) -> Result<PaginatedResponse<{}>>\n\
             - pub async fn get_by_id(pool, id) -> Result<{}>\n\
             - pub async fn insert(pool, payload) -> Result<{}>\n\
             - pub async fn update(pool, id, payload) -> Result<{}>\n\
             - pub async fn delete(pool, id) -> Result<bool>\n\
             {}\n\
             Use sqlx::query_as::<_, {}>() (runtime queries, not macros).\n\
             Use .bind() for parameters.\n",
            entity.name, entity.name, entity.name, entity.name,
            extra_queries_description, entity.name
        ),
        FileType::Service => format!(
            "This file must implement these business operations:\n\
             {}\n\
             Each function:\n\
             - Takes &PgPool, relevant params, &AuthUser\n\
             - Calls require_role() first\n\
             - Validates input\n\
             - Calls the appropriate query function\n\
             - Handles errors with AppError\n\
             - Logs with tracing::info!/warn!\n",
            business_rules_description
        ),
        // ... similar for API, Frontend, Tests
    }
}
\end{lstlisting}

% ============================================================
\section{Docker End-to-End}

The generated \texttt{docker-compose.yml} must actually start:

\begin{lstlisting}[title=Generated docker-compose.yml]
version: '3.8'
services:
  db:
    image: postgres:16-alpine
    environment:
      POSTGRES_DB: ${DB_NAME:-warehouse}
      POSTGRES_USER: ${DB_USER:-postgres}
      POSTGRES_PASSWORD: ${DB_PASSWORD:-postgres}
    volumes:
      - pgdata:/var/lib/postgresql/data
      - ./backend/database/migrations:/docker-entrypoint-initdb.d
    ports:
      - "5432:5432"
    healthcheck:
      test: ["CMD-SHELL", "pg_isready -U postgres"]
      interval: 5s
      timeout: 5s
      retries: 5

  backend:
    build: ./backend
    ports:
      - "${BACKEND_PORT:-8080}:8080"
    environment:
      DATABASE_URL: postgres://${DB_USER:-postgres}:${DB_PASSWORD:-postgres}@db/${DB_NAME:-warehouse}
      JWT_SECRET: ${JWT_SECRET:-dev-secret-change-in-production}
      RUST_LOG: info
    depends_on:
      db:
        condition: service_healthy

  frontend:
    image: nginx:alpine
    ports:
      - "${FRONTEND_PORT:-3000}:80"
    volumes:
      - ./frontend:/usr/share/nginx/html:ro
    depends_on:
      - backend

volumes:
  pgdata:
\end{lstlisting}

\subsection{Generated Dockerfile}

\begin{lstlisting}[title=Generated backend/Dockerfile]
FROM rust:1.77-slim AS builder
WORKDIR /app
COPY Cargo.toml Cargo.lock ./
COPY src ./src
RUN cargo build --release

FROM debian:bookworm-slim
RUN apt-get update && apt-get install -y libssl3 ca-certificates && rm -rf /var/lib/apt/lists/*
COPY --from=builder /app/target/release/{app_name} /usr/local/bin/server
EXPOSE 8080
CMD ["server"]
\end{lstlisting}

\subsection{SQL Migrations}

The LLM generates \texttt{database/migrations/001\_initial.sql}
with ALL tables, foreign keys, indices, and a seed admin user:

\begin{lstlisting}[language=SQL,title=End of 001_initial.sql]
-- Seed admin user (password: admin123, bcrypt hashed)
INSERT INTO users (email, password_hash, name, role, is_active, created_at, updated_at)
VALUES (
  'admin@example.com',
  '$2b$12$LJ3m4ys3Lz0Y8r5u.NXOxeVqH7VJvNRiKm4H1RU5q4v5iqN5rVEa',
  'Admin',
  'admin',
  true,
  NOW(),
  NOW()
) ON CONFLICT (email) DO NOTHING;
\end{lstlisting}

\subsection{Frontend CORS}

The backend must include CORS configuration that allows the
frontend (port 3000) to call the API (port 8080):

\begin{lstlisting}[language=C,title=In main.rs]
let cors = actix_cors::Cors::default()
    .allow_any_origin()
    .allow_any_method()
    .allow_any_header()
    .max_age(3600);
\end{lstlisting}

\subsection{End-to-End Verification}

\begin{lstlisting}[language=bash,title=The Test That Must Pass]
# Generate
isls forge-v2 --requirements examples/warehouse.toml \
    --api-key $OPENAI_API_KEY \
    --output ./output/warehouse

# Start
cd output/warehouse
docker-compose up -d

# Wait for startup
sleep 10

# Register a user
curl -s -X POST http://localhost:8080/api/auth/register \
  -H "Content-Type: application/json" \
  -d '{"email":"test@test.com","password":"test123","name":"Test User"}' \
  | jq .

# Login
TOKEN=$(curl -s -X POST http://localhost:8080/api/auth/login \
  -H "Content-Type: application/json" \
  -d '{"email":"admin@example.com","password":"admin123"}' \
  | jq -r '.token')

# Create product
curl -s -X POST http://localhost:8080/api/products \
  -H "Content-Type: application/json" \
  -H "Authorization: Bearer $TOKEN" \
  -d '{"sku":"WH-001","name":"Widget","category":"Parts",
       "unit_price_cents":1500,"quantity_on_hand":100,
       "reorder_level":10,"reorder_quantity":50,
       "warehouse_id":1}' \
  | jq .

# List products
curl -s http://localhost:8080/api/products \
  -H "Authorization: Bearer $TOKEN" | jq .

# Health check
curl -s http://localhost:8080/api/health | jq .

# Open browser
echo "Frontend: http://localhost:3000"
echo "API: http://localhost:8080/api/products"

# Cleanup
docker-compose down -v
\end{lstlisting}

% ============================================================
\section{Mock Mode}

With \texttt{--mock-oracle}, the system cannot call an LLM. In this
mode, use minimal hardcoded code for each file type:

\begin{itemize}
\item errors.rs: hardcoded AppError enum (always the same)
\item pagination.rs: hardcoded (always the same)
\item config.rs: hardcoded (always the same)
\item auth.rs: hardcoded (standard JWT implementation)
\item models: generate struct with all fields from norm (simple codegen,
  no LLM needed --- just format field names and types into a struct)
\item queries: generate from norm (format SQL strings with field names)
\item services: generate thin CRUD wrappers (validate + query + return)
\item api: generate thin handlers (extract params + call service + json)
\item main.rs: wire all routes
\item frontend: hardcoded SPA shell with API calls per entity
\end{itemize}

Mock mode produces thin but compilable code. LLM mode produces rich
code with real business logic. Both compile. Both start. LLM mode
is better.

Note: the ``mock generation'' functions are NOT Tera templates.
They are Rust functions that build code strings from norm artifacts:

\begin{lstlisting}[language=C]
fn mock_generate_model(entity: &EntityDef) -> String {
    let mut code = String::new();
    code.push_str("use serde::{Serialize, Deserialize};\n");
    code.push_str("use sqlx::FromRow;\n\n");
    // Main struct
    code.push_str(&format!(
        "#[derive(Debug, Clone, Serialize, Deserialize, FromRow)]\n\
         pub struct {} {{\n", entity.name));
    for field in &entity.fields {
        code.push_str(&format!("    pub {}: {},\n", field.name, field.rust_type));
    }
    code.push_str("}\n\n");
    // CreatePayload (only UserInput fields)
    code.push_str(&format!(
        "#[derive(Debug, Clone, Serialize, Deserialize)]\n\
         pub struct Create{} {{\n", entity.name));
    for field in &entity.fields {
        if matches!(field.source, FieldSource::UserInput) {
            code.push_str(&format!("    pub {}: {},\n", field.name, field.rust_type));
        }
    }
    code.push_str("}\n\n");
    // UpdatePayload (all fields as Option)
    // ... same pattern
    code
}
\end{lstlisting}

This is simple string formatting from norm data, not a template
engine. It cannot break because it reads the actual field names
and types from the norm.

% ============================================================
\section{Implementation Plan}

\begin{enumerate}
\item \textbf{Create isls-forge-llm module} in isls-renderloop (or new crate).
  Contains: LlmForge, TypeContext (growing), FileSpec, generation\_order().
  Static file generators (Cargo.toml, Docker, etc.).
  Mock generators (model, queries, service, api from norm data).
  Test: mock generates all files for warehouse, cargo check passes.

\item \textbf{LLM file generation + compile loop.}
  build\_prompt(), generate\_file(), cargo\_check(), retry with error.
  Test: generate errors.rs with LLM, cargo check passes.

\item \textbf{Growing type context.}
  add\_file() extracts types from generated file.
  Each subsequent prompt includes all previously generated types.
  Test: after generating errors.rs + models/product.rs, the prompt
  for product\_queries.rs includes both AppError and Product types.

\item \textbf{Full warehouse generation with LLM.}
  All layers generated in order. Each file compiles.
  Test: cargo build on complete backend succeeds.

\item \textbf{Docker end-to-end.}
  docker-compose.yml, Dockerfile, migrations with seed user, CORS.
  Test: docker-compose up, curl health check returns 200.

\item \textbf{Frontend generation.}
  Login page, entity pages, dashboard, API client.
  Test: nginx serves frontend, login form visible.

\item \textbf{End-to-end user flow.}
  Register, login, create product, list products, adjust stock.
  Test: the curl sequence from Section 4.5 succeeds.

\item \textbf{CLI integration.}
  \texttt{forge-v2 --api-key} uses LlmForge instead of templates.
  \texttt{forge-v2 --mock-oracle} uses mock generators.
  Test: both modes produce compilable output.
\end{enumerate}

% ============================================================
\section{Success Criteria}

\begin{enumerate}
\item \texttt{cargo build --workspace} zero errors.
\item All existing tests pass.
\item \texttt{forge-v2 --mock-oracle --output ./out/wh} $\to$
  backend compiles with 0 errors.
\item \texttt{forge-v2 --api-key \$KEY --output ./out/wh} $\to$
  backend compiles with 0 errors.
\item \texttt{docker-compose up -d} in generated output starts
  PostgreSQL + backend + frontend.
\item \texttt{curl /api/health} returns 200.
\item \texttt{curl /api/auth/login} with seed credentials returns JWT.
\item \texttt{curl /api/products} with JWT returns JSON (empty list or seeded data).
\item \texttt{curl -X POST /api/products} with JWT creates a product.
\item Frontend is accessible at \texttt{http://localhost:3000}.
\item LLM retries: if a file fails cargo check, the LLM gets the
  error and fixes it (logged in trace).
\item Generated code has NO Tera template artifacts (no \texttt{\{\{ \}\}},
  no \texttt{\{\% \%\}}).
\end{enumerate}

% ============================================================
\section{Constraints for Claude Code}

\begin{enumerate}
\item \textbf{Repository: graphity/isls.}
\item \textbf{Existing functionality must not break.}
  The old \texttt{forge-v2 --mock-oracle} path can use mock
  generators (not templates). The old template path is kept
  but no longer the default.
\item \textbf{OpenAI API calls must have timeouts.}
  30 second timeout per call. Use \texttt{reqwest::blocking}
  with timeout configured.
\item \textbf{cargo check runs in the generated backend directory.}
  Use \texttt{std::process::Command}. Parse stderr for errors.
\item \textbf{The growing type context is critical.}
  Every prompt must include ALL types generated so far.
  If the context gets too large ($>$ 10,000 tokens), summarize:
  show only struct names + field names, not full impls.
\item \textbf{Sequential generation.}
  Layer 0 before Layer 1 before Layer 2 before ... Layer 9.
  Within a layer: generate in entity dependency order.
\item \textbf{3 retry maximum per file.}
  If all 3 fail cargo check, log the error and continue to
  the next file. Do not abort the entire generation.
\item \textbf{Warehouse is the only target for v3.1.}
  Do not attempt ecommerce or project tracker. Get ONE domain
  perfect before generalizing.
\item \textbf{The docker-compose test is the most important
  deliverable.} If it starts and health check passes, the system
  works.
\item \textbf{No \texttt{unwrap()} in library code.}
\item \textbf{Every public item documented.}
\item \textbf{MIT License. Copyright (c) 2026 Sebastian Klemm.}
\end{enumerate}

\vfill
\noindent\rule{\textwidth}{0.4pt}
\begin{center}
\small
ISLS v3.1 --- Template-Free Generation --- Sebastian Klemm --- March 2026\\
Norms define WHAT. LLM generates HOW. Compiler verifies.\\
One domain. Warehouse. End-to-end. In the browser.\\
\url{https://github.com/LashSesh/graphity}
\end{center}

\end{document}
